\section{Conclusion}
\label{sec:conclusion}
The experiments clarify the roles of data generation, policy selection and program adaptation in robot self-training. Updating the data generator supports continued learning on fresh LIBERO scenes. In the BC experiments, selection preserves a capable policy and program adaptation enables further gains. Including RL's recorded tuning cost, the adapted loop achieves higher success on fresh LIBERO-Goal scenes with fewer GPU-hours than tuned one-shot RL.

\clearpage
\section*{LLM Usage Statement}
Large language models are part of the method studied here: the \texttt{llm} meta level (\cref{sec:meta}) calls Claude Sonnet 5.5 to propose edits to the training program; its inputs and outputs are logged and released with the run records. In the research and the writing we used Claude Code, an AI coding assistant (Claude Opus 5.5 and Claude Fable 5.1/5.2), to write and debug the experiment and analysis code, launch and monitor runs, compute the registered reads, generate the figures from the run records, and draft and edit text in all sections. The authors designed the study, fixed each registered claim and its analysis in a preregistration before its data existed, reviewed the code, numbers and text, and take full responsibility for the paper.
